\begin{proof}[Proof of \cref{obj:lem:global-holder-extension}]
Write
\[
\mathcal C:=[-1,1]^d.
\]
All intrinsic derivatives below are evaluated through their continuous traces, as in \cref{obj:def:holder}.

\begin{enumerate}
\item \emph{Order zero.}
Suppose \(m=0\). Define the coordinatewise retraction and the extension by
\[
\pi_{\mathcal C}:\mathbb R^d\longrightarrow\mathcal C,
\qquad
(\pi_{\mathcal C}(z))_i:=\max\{-1,\min\{1,z_i\}\},
\qquad
g_{\mathrm{ext}}(z):=g(\pi_{\mathcal C}(z)).
\]
Projection onto an interval is nonexpansive. Consequently,
\[
\pi_{\mathcal C}(z)=z\quad(z\in\mathcal C),
\qquad
\|\pi_{\mathcal C}(z)-\pi_{\mathcal C}(z')\|_\infty
\le \|z-z'\|_\infty.
\]
The hypotheses therefore give
\[
|g_{\mathrm{ext}}(z)|\le L,
\qquad
|g_{\mathrm{ext}}(z)-g_{\mathrm{ext}}(z')|
\le L\|\pi_{\mathcal C}(z)-\pi_{\mathcal C}(z')\|_\infty^s
\le L\|z-z'\|_\infty^s.
\]
Since \(s>0\), the last bound implies continuity. Thus \(K_{\mathrm{ext}}=1\) works in this case. The same construction applies when \(d=0\), with the empty coordinate tuple interpreted as the unique point of \(\mathbb R^0\).
% lean: Causalean.Mathlib.Analysis.Calculus.CubeExtension.exists_global_holder_extension_order_zero

\item\label{proof:global-holder-extension:coordinate-bounds} \emph{From coordinate bounds to multilinear bounds.}
Henceforth assume \(m>0\). Set
\[
K_{\mathrm{coord}}:=(d+1)^m>0.
\]
For endpoints \(\boldsymbol\ell,\boldsymbol b\in\mathbb R^d\) satisfying
\(\ell_i<b_i\), define
\[
\mathcal B(\boldsymbol\ell,\boldsymbol b)
:=\prod_{i=1}^d[\ell_i,b_i].
\]
For an intrinsic \(C^m\) response \(g_{\mathrm{in}}\) on this box, its full intrinsic derivative is the multilinear form
\[
\mathrm D_{\mathcal B}^{k}g_{\mathrm{in}}(z)
[\xi_1,\ldots,\xi_k]
:=
\sum_{i_1,\ldots,i_k=1}^d
\bigl(\partial_{i_1}\cdots\partial_{i_k}g_{\mathrm{in}}\bigr)(z)
\prod_{\nu=1}^k(\xi_\nu)_{i_\nu},
\qquad 0\le k\le m.
\]
At order zero this is \(g_{\mathrm{in}}(z)\); at positive order in dimension zero the sum is empty and the form is zero.

To see the relevant norm estimate, define the coordinate basis by
\[
(\mathbf e_i)_j:=\mathbf 1_{\{i=j\}}.
\]
If a \(k\)-linear form \(M_{\mathrm{coord}}:(\mathbb R^d)^k\to\mathbb R\) satisfies
\[
|M_{\mathrm{coord}}[\mathbf e_{i_1},\ldots,\mathbf e_{i_k}]|
\le L_{\mathrm{coord}},
\qquad L_{\mathrm{coord}}\ge0,
\]
multilinearity gives
\[
\begin{aligned}
|M_{\mathrm{coord}}[\xi_1,\ldots,\xi_k]|
&\le
\sum_{i_1,\ldots,i_k=1}^d
L_{\mathrm{coord}}\prod_{\nu=1}^k|(\xi_\nu)_{i_\nu}|\\
&\le d^kL_{\mathrm{coord}}
\prod_{\nu=1}^k\|\xi_\nu\|_\infty
\le K_{\mathrm{coord}}L_{\mathrm{coord}}
\prod_{\nu=1}^k\|\xi_\nu\|_\infty.
\end{aligned}
\]
Here the order-zero sum has one term, including when \(d=0\). Applying this estimate to the derivative and to the difference of two top derivatives shows that coordinate bounds of radius \(L_{\mathrm{in}}\ge0\) imply
\[
\|\mathrm D_{\mathcal B}^{k}g_{\mathrm{in}}(z)\|_{\mathrm{op}}
\le K_{\mathrm{coord}}L_{\mathrm{in}},
\qquad 0\le k\le m,\quad z\in\mathcal B,
\]
and
\[
\|\mathrm D_{\mathcal B}^{m}g_{\mathrm{in}}(z)
-\mathrm D_{\mathcal B}^{m}g_{\mathrm{in}}(z')\|_{\mathrm{op}}
\le K_{\mathrm{coord}}L_{\mathrm{in}}\|z-z'\|_\infty^s,
\qquad z,z'\in\mathcal B.
\]
% lean: Causalean.Mathlib.Analysis.Calculus.CubeExtension.exists_coord_multilinear_norm_constant

\item\label{proof:global-holder-extension:affine-transport} \emph{Affine transport between boxes.}
Consider two nondegenerate boxes
\[
\mathcal B=\mathcal B(\boldsymbol\ell,\boldsymbol b),
\qquad
\mathcal B'=\mathcal B(\boldsymbol\ell',\boldsymbol b').
\]
Define their affine identification and its coordinate slopes by
\[
\lambda_{\mathrm{aff},i}
:=\frac{b_i'-\ell_i'}{b_i-\ell_i}>0,
\qquad
(\Phi_{\mathcal B\to\mathcal B'}(z))_i
:=\ell_i'+\lambda_{\mathrm{aff},i}(z_i-\ell_i).
\]
This map carries \(\mathcal B\) bijectively onto \(\mathcal B'\). Define
\[
K_{\mathrm{aff},\mathrm{jet}}
:=1+\sum_{i=1}^d\lambda_{\mathrm{aff},i},
\qquad
K_{\mathrm{aff},\mathrm{dist}}
:=1+\sum_{i=1}^d
\bigl(\lambda_{\mathrm{aff},i}+\lambda_{\mathrm{aff},i}^{-1}\bigr),
\]
and
\[
K_{\mathrm{aff}}(\mathcal B,\mathcal B')
:=
K_{\mathrm{aff},\mathrm{jet}}^m
\bigl(1+K_{\mathrm{aff},\mathrm{dist}}^s\bigr)>0.
\]
Coordinatewise subtraction yields
\[
\|\Phi_{\mathcal B\to\mathcal B'}(z)
-\Phi_{\mathcal B\to\mathcal B'}(z')\|_\infty
\le K_{\mathrm{aff},\mathrm{dist}}\|z-z'\|_\infty,
\]
hence, because \(s>0\),
\[
\|\Phi_{\mathcal B\to\mathcal B'}(z)
-\Phi_{\mathcal B\to\mathcal B'}(z')\|_\infty^s
\le K_{\mathrm{aff},\mathrm{dist}}^s\|z-z'\|_\infty^s.
\]
For a response \(g_{\mathrm{in}}\) with intrinsic radius \(L_{\mathrm{in}}\) on \(\mathcal B'\), define
\[
g_{\mathrm{aff}}:=g_{\mathrm{in}}\circ\Phi_{\mathcal B\to\mathcal B'}.
\]
The affine chain rule, including continuous boundary traces, gives
\[
D^\alpha g_{\mathrm{aff}}(z)
=
\left(\prod_{i=1}^d\lambda_{\mathrm{aff},i}^{\alpha_i}\right)
(D^\alpha g_{\mathrm{in}})
(\Phi_{\mathcal B\to\mathcal B'}(z)).
\]
Since \(K_{\mathrm{aff},\mathrm{jet}}\ge1\), for \(|\alpha|\le m\),
\[
\prod_{i=1}^d\lambda_{\mathrm{aff},i}^{\alpha_i}
\le K_{\mathrm{aff},\mathrm{jet}}^{|\alpha|}
\le K_{\mathrm{aff},\mathrm{jet}}^m.
\]
It follows that
\[
|D^\alpha g_{\mathrm{aff}}(z)|
\le K_{\mathrm{aff},\mathrm{jet}}^mL_{\mathrm{in}}
\le K_{\mathrm{aff}}(\mathcal B,\mathcal B')L_{\mathrm{in}},
\]
and, for \(|\alpha|=m\),
\[
\begin{aligned}
|D^\alpha g_{\mathrm{aff}}(z)-D^\alpha g_{\mathrm{aff}}(z')|
&\le K_{\mathrm{aff},\mathrm{jet}}^mL_{\mathrm{in}}
\|\Phi_{\mathcal B\to\mathcal B'}(z)
-\Phi_{\mathcal B\to\mathcal B'}(z')\|_\infty^s\\
&\le K_{\mathrm{aff}}(\mathcal B,\mathcal B')L_{\mathrm{in}}
\|z-z'\|_\infty^s.
\end{aligned}
\]
Composition with the affine map also preserves intrinsic \(C^m\) regularity. These formulas establish the claimed radius factor over the whole closed box.
% lean: Causalean.Mathlib.Analysis.Calculus.CubeExtension.exists_rectAffine_holder_constant

\item \emph{Moment-matched reflection across a normalized face.}
Choose real coefficients \(a_{\mathrm{ref},r}\), \(1\le r\le m+1\), satisfying
\[
\sum_{r=1}^{m+1}a_{\mathrm{ref},r}(-r)^k=1,
\qquad 0\le k\le m.
\]
Such coefficients exist: the matrix
\[
(\mathsf V_{\mathrm{ref}})_{r,k+1}:=(-r)^k,
\qquad 1\le r\le m+1,\quad 0\le k\le m,
\]
is a Vandermonde matrix at distinct nodes. Its determinant is nonzero, so the system
\[
\mathsf V_{\mathrm{ref}}^{\mathsf T}
(a_{\mathrm{ref},1},\ldots,a_{\mathrm{ref},m+1})^{\mathsf T}
=(1,\ldots,1)^{\mathsf T}
\]
has a solution depending only on \(m\).
% lean: Causalean.Mathlib.Analysis.Calculus.CubeExtension.exists_reflection_coefficients

\item\label{proof:global-holder-extension:face-collar} \emph{Regularity and bounds on a one-face collar.}
Fix a coordinate \(i\), and define
\[
\mathcal C_i^-:=
\left\{z:
-1-\frac1{m+1}\le z_i\le1,\
-1\le z_j\le1\ \text{for }j\ne i
\right\},
\]
\[
\mathcal E_i^-:=\mathcal C_i^-\cap\{z:z_i\le-1\},
\qquad
\mathcal E_i^{-,\circ}:=
\left\{z:
-1-\frac1{m+1}<z_i<-1,\
-1<z_j<1\ \text{for }j\ne i
\right\}.
\]
For each \(1\le r\le m+1\), define an affine inward sample on all of \(\mathbb R^d\) by
\[
(\Phi_{\mathrm{ref},i,r}(z))_j
:=
\begin{cases}
-1+r(-1-z_i),&j=i,\\
z_j,&j\ne i.
\end{cases}
\]
If \(z\in\mathcal E_i^{-,\circ}\), then
\[
0<r(-1-z_i)<1,
\]
so \(\Phi_{\mathrm{ref},i,r}(z)\in(-1,1)^d\). The weak inequalities likewise show that the closed exterior slab maps into \(\mathcal C\).

For an auxiliary response \(g_{\mathrm{in}}:\mathbb R^d\to\mathbb R\) with intrinsic radius \(L_{\mathrm{in}}\) on \(\mathcal C\), define
\[
(\mathscr R_i g_{\mathrm{in}})(z)
:=
\begin{cases}
\displaystyle\sum_{r=1}^{m+1}
a_{\mathrm{ref},r}g_{\mathrm{in}}(\Phi_{\mathrm{ref},i,r}(z)),
&z_i<-1,\\
g_{\mathrm{in}}(z),&z_i\ge-1.
\end{cases}
\]
This defines the response on all of \(\mathbb R^d\), and it agrees with \(g_{\mathrm{in}}\) on \(\mathcal C\).

On the closed exterior slab, the finite sum is intrinsically \(C^m\), by composition with the affine samples. At \(z_i=-1\), every sample equals \(z\), and the moment equation at \(k=0\) identifies that sum with the second branch. Moreover, for \(|\alpha|\le m\), the exterior derivative trace at this face is
\[
\begin{aligned}
(D^\alpha\mathscr R_i g_{\mathrm{in}})_{\mathrm{ext}}(z)
&=
\sum_{r=1}^{m+1}
a_{\mathrm{ref},r}(-r)^{\alpha_i}D^\alpha g_{\mathrm{in}}(z)\\
&=D^\alpha g_{\mathrm{in}}(z).
\end{aligned}
\]
Indeed, the affine chain rule gives the factors \((-r)^{\alpha_i}\), and continuity of the original derivative traces gives the limit at the face. Coordinate expansion then identifies the full multilinear traces. This argument includes the edges and corners.

The two closed pieces therefore have matching derivatives through order \(m\). They glue to an intrinsic \(C^m\) response on \(\mathcal C_i^-\): at a common face point, the derivative remainders on both pieces use the same linear approximation and are \(o(\|\xi\|_\infty)\); on their union the same remainder estimate holds. Applying this observation successively to the derivative fields, and using continuity of the matching fields, establishes regularity through order \(m\).

For the quantitative estimates, define the linear part of each sample by
\[
(\Lambda_{\mathrm{ref},i,r}\xi)_j
:=
\begin{cases}
-r\xi_i,&j=i,\\
\xi_j,&j\ne i.
\end{cases}
\]
Then
\[
\|\Lambda_{\mathrm{ref},i,r}\|_{\mathrm{op}}\le r,
\qquad
\|\Phi_{\mathrm{ref},i,r}(z)-\Phi_{\mathrm{ref},i,r}(z')\|_\infty
\le r\|z-z'\|_\infty.
\]
On the open exterior slab, the chain rule gives
\[
\mathrm D^k(g_{\mathrm{in}}\circ\Phi_{\mathrm{ref},i,r})(z)
=
\mathrm D^k g_{\mathrm{in}}(\Phi_{\mathrm{ref},i,r}(z))
\circ\Lambda_{\mathrm{ref},i,r}^{\otimes k}.
\]
Using the full intrinsic bounds from \cref{proof:global-holder-extension:coordinate-bounds} at the inward samples yields
\[
\|\mathrm D^k(g_{\mathrm{in}}\circ\Phi_{\mathrm{ref},i,r})(z)\|_{\mathrm{op}}
\le K_{\mathrm{coord}}L_{\mathrm{in}}r^k
\le K_{\mathrm{coord}}r^mL_{\mathrm{in}},
\]
and
\[
\begin{aligned}
&\|\mathrm D^m(g_{\mathrm{in}}\circ\Phi_{\mathrm{ref},i,r})(z)
-\mathrm D^m(g_{\mathrm{in}}\circ\Phi_{\mathrm{ref},i,r})(z')\|_{\mathrm{op}}\\
&\qquad\le
K_{\mathrm{coord}}L_{\mathrm{in}}r^m
\|\Phi_{\mathrm{ref},i,r}(z)-\Phi_{\mathrm{ref},i,r}(z')\|_\infty^s\\
&\qquad\le K_{\mathrm{coord}}r^mr^sL_{\mathrm{in}}\|z-z'\|_\infty^s.
\end{aligned}
\]
Define
\[
K_{\mathrm{ref},\mathrm{ext},D}
:=1+\sum_{r=1}^{m+1}|a_{\mathrm{ref},r}|K_{\mathrm{coord}}r^m,
\qquad
K_{\mathrm{ref},\mathrm{ext},H}
:=1+\sum_{r=1}^{m+1}|a_{\mathrm{ref},r}|K_{\mathrm{coord}}r^mr^s.
\]
Differentiating the finite reflection sum and applying the triangle inequality gives
\[
\|\mathrm D^k\mathscr R_i g_{\mathrm{in}}(z)\|_{\mathrm{op}}
\le K_{\mathrm{ref},\mathrm{ext},D}L_{\mathrm{in}},
\]
and
\[
\|\mathrm D^m\mathscr R_i g_{\mathrm{in}}(z)
-\mathrm D^m\mathscr R_i g_{\mathrm{in}}(z')\|_{\mathrm{op}}
\le K_{\mathrm{ref},\mathrm{ext},H}L_{\mathrm{in}}\|z-z'\|_\infty^s
\]
on the open exterior slab. Its closure is \(\mathcal E_i^-\). Continuity of the within-collar derivatives therefore extends the first bound to that closure, and successive limits in the two points extend the second bound.

On \(\mathcal C\), agreement of the responses identifies their derivatives in the interior. Continuity extends this identity to the closed cube:
\[
\mathrm D_{\mathcal C_i^-}^k\mathscr R_i g_{\mathrm{in}}(z)
=\mathrm D_{\mathcal C}^k g_{\mathrm{in}}(z),
\qquad z\in\mathcal C,\quad k\le m.
\]
Thus the corresponding full derivative and top-modulus bounds on the cube have coefficient \(K_{\mathrm{coord}}\).

For a pair \(z\in\mathcal E_i^-\), \(z'\in\mathcal C\), choose
\[
t_{\mathrm{face}}\in[0,1],
\qquad
(1-t_{\mathrm{face}})z_i+t_{\mathrm{face}}z_i'=-1,
\qquad
z_{\mathrm{face}}:=(1-t_{\mathrm{face}})z+t_{\mathrm{face}}z'.
\]
Existence follows from \(z_i\le-1\le z_i'\); if both coordinates equal \(-1\), take \(t_{\mathrm{face}}=0\). The other coordinates remain in \([-1,1]\), so
\[
z_{\mathrm{face}}\in\mathcal E_i^-\cap\mathcal C,
\qquad
\|z-z_{\mathrm{face}}\|_\infty\le\|z-z'\|_\infty,
\qquad
\|z_{\mathrm{face}}-z'\|_\infty\le\|z-z'\|_\infty.
\]
The two side estimates and the triangle inequality now give
\[
\begin{aligned}
&\|\mathrm D_{\mathcal C_i^-}^m\mathscr R_i g_{\mathrm{in}}(z)
-\mathrm D_{\mathcal C_i^-}^m\mathscr R_i g_{\mathrm{in}}(z')\|_{\mathrm{op}}\\
&\quad\le
K_{\mathrm{ref},\mathrm{ext},H}L_{\mathrm{in}}
\|z-z_{\mathrm{face}}\|_\infty^s
+K_{\mathrm{coord}}L_{\mathrm{in}}
\|z_{\mathrm{face}}-z'\|_\infty^s\\
&\quad\le
(K_{\mathrm{ref},\mathrm{ext},H}+K_{\mathrm{coord}})
L_{\mathrm{in}}\|z-z'\|_\infty^s.
\end{aligned}
\]
Pairs on the same side satisfy the same enlarged bound; reversing the points treats the other crossing order. Evaluating on coordinate basis tuples converts these full-jet estimates to coordinate estimates. Consequently, with
\[
K_{\mathrm{ref},D}
:=K_{\mathrm{ref},\mathrm{ext},D}+K_{\mathrm{coord}},
\qquad
K_{\mathrm{ref},H}
:=K_{\mathrm{ref},\mathrm{ext},H}+K_{\mathrm{coord}},
\qquad
K_{\mathrm{ref}}:=K_{\mathrm{ref},D}+K_{\mathrm{ref},H},
\]
the reflected response has intrinsic \(C^{m,s}\) radius
\(K_{\mathrm{ref}}L_{\mathrm{in}}\) on \(\mathcal C_i^-\). All these constants are positive and depend only on \(d,m,s\).
% lean: Causalean.Mathlib.Analysis.Calculus.CubeExtension.exists_leftFaceReflection_collar_holder_constant

\item\label{proof:global-holder-extension:larger-box} \emph{A fixed larger box.}
For a nondegenerate box
\(\mathcal B=\mathcal B(\boldsymbol\ell,\boldsymbol b)\) and a coordinate \(i\), define its lower enlargement by
\[
\ell_j^-:=
\begin{cases}
\displaystyle\ell_i-\frac{b_i-\ell_i}{2(m+1)},&j=i,\\
\ell_j,&j\ne i,
\end{cases}
\qquad
\mathcal B^-:=\mathcal B(\boldsymbol\ell^-,\boldsymbol b).
\]
Normalize a response on \(\mathcal B\), reflect it, and transport it back:
\[
g_{\mathrm{norm}}
:=g_{\mathrm{in}}\circ\Phi_{\mathcal C\to\mathcal B},
\qquad
g_{\mathrm{lower}}
:=(\mathscr R_i g_{\mathrm{norm}})
\circ\Phi_{\mathcal B^-\to\mathcal C_i^-}.
\]
The chosen collar width gives the affine identity
\[
\Phi_{\mathcal B^-\to\mathcal C_i^-}
=\Phi_{\mathcal B\to\mathcal C}
\quad\text{on }\mathbb R^d.
\]
Therefore
\[
g_{\mathrm{lower}}(z)=g_{\mathrm{in}}(z)
\quad(z\in\mathcal B).
\]
\Cref{proof:global-holder-extension:affine-transport,proof:global-holder-extension:face-collar} give its radius factor
\[
K_{\mathrm{lower}}(\mathcal B,i)
:=
K_{\mathrm{aff}}(\mathcal B^-,\mathcal C_i^-)
K_{\mathrm{ref}}
K_{\mathrm{aff}}(\mathcal C,\mathcal B)>0.
\]

For the upper enlargement, define
\[
b_j^+:=
\begin{cases}
\displaystyle b_i+\frac{b_i-\ell_i}{2(m+1)},&j=i,\\
b_j,&j\ne i,
\end{cases}
\qquad
\mathcal B^+:=\mathcal B(\boldsymbol\ell,\boldsymbol b^+).
\]
Apply the lower construction to the negated box and then negate back. This transport preserves the radius, since
\[
D^\alpha(g_{\mathrm{in}}\circ(-\mathrm{id}))(z)
=(-1)^{|\alpha|}D^\alpha g_{\mathrm{in}}(-z),
\qquad
\|(-z)-(-z')\|_\infty=\|z-z'\|_\infty.
\]
Thus the upper construction agrees on \(\mathcal B\) and has factor
\[
K_{\mathrm{upper}}(\mathcal B,i)
:=K_{\mathrm{lower}}(-\mathcal B,i),
\qquad
-\mathcal B:=\mathcal B(-\boldsymbol b,-\boldsymbol\ell).
\]
Applying the lower construction first and the upper construction to its enlarged box gives agreement on the original box and factor
\[
K_{\mathrm{two}}(\mathcal B,i)
:=K_{\mathrm{upper}}(\mathcal B^-,i)
K_{\mathrm{lower}}(\mathcal B,i)>0.
\]

Iterate this two-face construction through coordinates \(1,\ldots,d\). More precisely, define
\[
\boldsymbol\ell^{(0)}:=(-1,\ldots,-1),
\qquad
\boldsymbol b^{(0)}:=(1,\ldots,1),
\qquad
\mathcal B^{(0)}:=\mathcal C,
\qquad
K_{\mathrm{neigh}}^{(0)}:=1.
\]
At stage \(\nu=1,\ldots,d\), apply the lower enlargement in coordinate \(\nu\), followed by the upper enlargement using the new lower endpoint, and denote the resulting box by \(\mathcal B^{(\nu)}\). Define
\[
K_{\mathrm{neigh}}^{(\nu)}
:=
K_{\mathrm{two}}(\mathcal B^{(\nu-1)},\nu)
K_{\mathrm{neigh}}^{(\nu-1)}.
\]
The lower and upper increments are positive. Previously enlarged coordinates remain unchanged, and every intermediate box contains \(\mathcal C\). Induction consequently gives
\[
\ell_i^{(d)}<-1<1<b_i^{(d)}
\quad(1\le i\le d),
\qquad
K_{\mathrm{neigh}}^{(d)}>0.
\]
Set
\[
\mathcal B_{\mathrm{ext}}
:=\mathcal B^{(d)},
\qquad
\boldsymbol\ell_{\mathrm{ext}}:=\boldsymbol\ell^{(d)},
\qquad
\boldsymbol b_{\mathrm{ext}}:=\boldsymbol b^{(d)},
\qquad
K_{\mathrm{neigh}}:=K_{\mathrm{neigh}}^{(d)}.
\]
For every admissible \(g,L\), the same induction produces a response \(g_{\mathrm{neigh}}\) satisfying
\[
g_{\mathrm{neigh}}=g\quad\text{on }\mathcal C,
\qquad
g_{\mathrm{neigh}}\in C^m(\mathcal B_{\mathrm{ext}}),
\]
\[
|D^\alpha g_{\mathrm{neigh}}(z)|
\le K_{\mathrm{neigh}}L
\quad(|\alpha|\le m,\ z\in\mathcal B_{\mathrm{ext}}),
\]
and
\[
|D^\alpha g_{\mathrm{neigh}}(z)-D^\alpha g_{\mathrm{neigh}}(z')|
\le K_{\mathrm{neigh}}L\|z-z'\|_\infty^s
\quad(|\alpha|=m,\ z,z'\in\mathcal B_{\mathrm{ext}}).
\]
For definiteness, the starting ambient representative is
\[
g^{(0)}(z):=
\begin{cases}
g(z),&z\in\mathcal C,\\
0,&z\notin\mathcal C.
\end{cases}
\]
Subsequent representatives are defined everywhere by the displayed reflection and affine formulas. When \(d=0\), the iteration is empty:
\(\mathcal B_{\mathrm{ext}}=\mathbb R^0\) and \(K_{\mathrm{neigh}}=1\).
The final box and radius factor depend only on \(d,m,s\).
% lean: Causalean.Mathlib.Analysis.Calculus.CubeExtension.exists_fixedCubeNeighborhood_holder_constant

\item\label{proof:global-holder-extension:cutoff} \emph{A fixed smooth cutoff and its zero extension.}
Define
\[
\rho_{\mathrm{box}}:=
\begin{cases}
\displaystyle\min_{1\le i\le d}
\min\{-\ell_{\mathrm{ext},i},b_{\mathrm{ext},i}\},&d>0,\\
2,&d=0,
\end{cases}
\]
and
\[
\rho_{\mathrm{in}}:=\frac{1+\rho_{\mathrm{box}}}{2},
\qquad
\rho_{\mathrm{out}}:=\frac{\rho_{\mathrm{in}}+\rho_{\mathrm{box}}}{2}.
\]
The strict margins imply
\[
1<\rho_{\mathrm{in}}<\rho_{\mathrm{out}}<\rho_{\mathrm{box}}.
\]
Choose a smooth bump \(\chi:\mathbb R^d\to\mathbb R\) with
\[
\chi(z)=1\quad(\|z\|_\infty\le\rho_{\mathrm{in}}),
\qquad
\operatorname{supp}\chi
\subseteq\{z:\|z\|_\infty\le\rho_{\mathrm{out}}\},
\]
where support denotes closed support. Smooth bump existence in a finite-dimensional normed space supplies this choice. In particular,
\[
\chi=1\quad\text{on }\mathcal C,
\qquad
\operatorname{supp}\chi\subset
\mathcal B_{\mathrm{ext}}^\circ,
\]
with
\[
\mathcal B_{\mathrm{ext}}^\circ
:=\prod_{i=1}^d
(\ell_{\mathrm{ext},i},b_{\mathrm{ext},i}).
\]
Compact support and continuity of the derivatives make the constant
\[
K_\chi:=
\max\left\{
1,\ 
\max_{0\le k\le m+1}
\sup_{z\in\mathbb R^d}
\|\mathrm D^k\chi(z)\|_{\mathrm{op}}
\right\}
\]
finite and positive.

For any response \(g_{\mathrm{in}}\) of intrinsic radius \(L_{\mathrm{in}}\ge0\) on \(\mathcal B_{\mathrm{ext}}\), define
\[
(\mathscr G g_{\mathrm{in}})(z)
:=
\begin{cases}
\chi(z)g_{\mathrm{in}}(z),&z\in\mathcal B_{\mathrm{ext}},\\
0,&z\notin\mathcal B_{\mathrm{ext}}.
\end{cases}
\]
At an interior point this is locally a product of \(C^m\) functions. At every point outside the open box, a neighborhood misses the closed support of \(\chi\), so the extension is locally zero. Hence
\[
\mathscr G g_{\mathrm{in}}\in C^m(\mathbb R^d),
\qquad
\mathscr G g_{\mathrm{in}}=g_{\mathrm{in}}\quad\text{on }\mathcal C,
\]
and
\[
\mathrm D^k\mathscr G g_{\mathrm{in}}(z)=0
\quad(z\notin\mathcal B_{\mathrm{ext}}^\circ,\ 0\le k\le m).
\]
% lean: Causalean.Mathlib.Analysis.Calculus.CubeExtension.rectCutoffExtension_contDiff_eqOn_cube

\item\label{proof:global-holder-extension:cutoff-derivatives} \emph{Uniform derivatives of the cutoff extension.}
We make the product estimate explicit. Define
\[
\mathfrak S_k:=\{\text{permutations of }\{1,\ldots,k\}\}.
\]
For an \(r\)-linear form \(M_1\) and a \((k-r)\)-linear form \(M_2\), where \(0\le r\le k\), define
\[
\begin{aligned}
&\mathcal M_{k,r}(M_1,M_2)[\xi_1,\ldots,\xi_k]\\
&\qquad:=
\frac1{k!}\sum_{\sigma\in\mathfrak S_k}
M_1[\xi_{\sigma(1)},\ldots,\xi_{\sigma(r)}]\,
M_2[\xi_{\sigma(r+1)},\ldots,\xi_{\sigma(k)}].
\end{aligned}
\]
At \(k=0\), the permutation set has one element and the expression is ordinary multiplication. Each summand is bounded by
\(\|M_1\|_{\mathrm{op}}\|M_2\|_{\mathrm{op}}
\prod_{\nu=1}^k\|\xi_\nu\|_\infty\), so
\[
\|\mathcal M_{k,r}(M_1,M_2)\|_{\mathrm{op}}
\le\|M_1\|_{\mathrm{op}}\|M_2\|_{\mathrm{op}}.
\]
Repeated differentiation along a line gives the scalar Leibniz formula on repeated arguments. Both sides below are symmetric multilinear forms, and polarization therefore gives the full product identity
\[
\mathrm D^k(\chi g_{\mathrm{in}})(z)
=
\sum_{r=0}^k\binom{k}{r}
\mathcal M_{k,r}
\bigl(\mathrm D^r\chi(z),
\mathrm D_{\mathcal B_{\mathrm{ext}}}^{k-r}g_{\mathrm{in}}(z)\bigr),
\qquad z\in\mathcal B_{\mathrm{ext}}^\circ.
\]
Thus \cref{proof:global-holder-extension:coordinate-bounds} gives
\[
\begin{aligned}
\|\mathrm D^k\mathscr G g_{\mathrm{in}}(z)\|_{\mathrm{op}}
&\le
\sum_{r=0}^k\binom{k}{r}
K_\chi K_{\mathrm{coord}}L_{\mathrm{in}}\\
&=2^kK_\chi K_{\mathrm{coord}}L_{\mathrm{in}}
\le2^mK_\chi K_{\mathrm{coord}}L_{\mathrm{in}}.
\end{aligned}
\]
Outside the open box the derivative is zero. Consequently, defining
\[
K_{\mathrm{cut},D}:=K_\chi K_{\mathrm{coord}}2^m+1>0,
\]
we obtain the global bound
\[
\|\mathrm D^k\mathscr G g_{\mathrm{in}}(z)\|_{\mathrm{op}}
\le K_{\mathrm{cut},D}L_{\mathrm{in}},
\qquad 0\le k\le m,\quad z\in\mathbb R^d.
\]
% lean: Causalean.Mathlib.Analysis.Calculus.CubeExtension.exists_rectCutoffExtension_jet_constant

\item\label{proof:global-holder-extension:global-modulus} \emph{The global top-order modulus.}
For \(z,z'\) in the open box, define
\[
\delta:=\|z-z'\|_\infty.
\]
The box is convex. For \(k<m\), the bound on the next full intrinsic derivative and the mean-value inequality give
\[
\|\mathrm D_{\mathcal B_{\mathrm{ext}}}^{k}g_{\mathrm{in}}(z)
-\mathrm D_{\mathcal B_{\mathrm{ext}}}^{k}g_{\mathrm{in}}(z')\|_{\mathrm{op}}
\le K_{\mathrm{coord}}L_{\mathrm{in}}\delta.
\]
The uniform bounds also give a difference bound
\(2K_{\mathrm{coord}}L_{\mathrm{in}}\). For \(0\le\delta\le1\), the condition \(s\le1\) implies \(\delta\le\delta^s\); for \(\delta>1\), positivity of \(s\) implies \(1\le\delta^s\). Therefore
\[
\|\mathrm D_{\mathcal B_{\mathrm{ext}}}^{k}g_{\mathrm{in}}(z)
-\mathrm D_{\mathcal B_{\mathrm{ext}}}^{k}g_{\mathrm{in}}(z')\|_{\mathrm{op}}
\le2K_{\mathrm{coord}}L_{\mathrm{in}}\delta^s,
\qquad k<m.
\]
For \(k=m\), \cref{proof:global-holder-extension:coordinate-bounds} gives the sharper bound with coefficient
\(K_{\mathrm{coord}}L_{\mathrm{in}}\), and hence the same displayed bound with coefficient \(2K_{\mathrm{coord}}L_{\mathrm{in}}\).

Likewise, the cutoff derivatives through order \(m+1\) give, for \(0\le r\le m\),
\[
\|\mathrm D^r\chi(z)-\mathrm D^r\chi(z')\|_{\mathrm{op}}
\le K_\chi\delta,
\qquad
\|\mathrm D^r\chi(z)-\mathrm D^r\chi(z')\|_{\mathrm{op}}
\le2K_\chi.
\]
The same two distance ranges yield
\[
\|\mathrm D^r\chi(z)-\mathrm D^r\chi(z')\|_{\mathrm{op}}
\le2K_\chi\delta^s.
\]

Bilinearity of the pairings in \cref{proof:global-holder-extension:cutoff-derivatives} splits each product difference into a cutoff difference times a response jet, and a cutoff jet times a response difference. Subtracting the two product identities and applying their contraction bounds gives
\[
\begin{aligned}
&\|\mathrm D^m\mathscr G g_{\mathrm{in}}(z)
-\mathrm D^m\mathscr G g_{\mathrm{in}}(z')\|_{\mathrm{op}}\\
&\quad\le
\sum_{r=0}^m\binom mr
\Bigl(
\|\mathrm D^r\chi(z)-\mathrm D^r\chi(z')\|_{\mathrm{op}}
\|\mathrm D_{\mathcal B_{\mathrm{ext}}}^{m-r}g_{\mathrm{in}}(z)\|_{\mathrm{op}}\\
&\hspace{105pt}
+\|\mathrm D^r\chi(z')\|_{\mathrm{op}}
\|\mathrm D_{\mathcal B_{\mathrm{ext}}}^{m-r}g_{\mathrm{in}}(z)
-\mathrm D_{\mathcal B_{\mathrm{ext}}}^{m-r}g_{\mathrm{in}}(z')\|_{\mathrm{op}}
\Bigr)\\
&\quad\le
\sum_{r=0}^m\binom mr
\left(
2K_\chi\delta^s K_{\mathrm{coord}}L_{\mathrm{in}}
+K_\chi\,2K_{\mathrm{coord}}L_{\mathrm{in}}\delta^s
\right)\\
&\quad=
4K_\chi K_{\mathrm{coord}}2^mL_{\mathrm{in}}\delta^s.
\end{aligned}
\]
Define
\[
K_{\mathrm{cut},H}:=4K_\chi K_{\mathrm{coord}}2^m+1>0.
\]
The open box is dense in the closed box, including in dimension zero. Since \(\mathscr G g_{\mathrm{in}}\) is globally \(C^m\), its top derivative is continuous. Passing to limits in both points, using continuity of the distance power for \(s>0\), gives
\[
\|\mathrm D^m\mathscr G g_{\mathrm{in}}(z)
-\mathrm D^m\mathscr G g_{\mathrm{in}}(z')\|_{\mathrm{op}}
\le K_{\mathrm{cut},H}L_{\mathrm{in}}\|z-z'\|_\infty^s,
\qquad z,z'\in\mathcal B_{\mathrm{ext}}.
\]

Now let \(z\in\mathcal B_{\mathrm{ext}}\) and \(z'\notin\mathcal B_{\mathrm{ext}}\). Some coordinate \(i\) satisfies
\(z_i'<\ell_{\mathrm{ext},i}\) or \(z_i'>b_{\mathrm{ext},i}\). Define
\[
(z_{\mathrm{boundary}})_j:=
\begin{cases}
\ell_{\mathrm{ext},i},&j=i,\ z_i'<\ell_{\mathrm{ext},i},\\
b_{\mathrm{ext},i},&j=i,\ z_i'>b_{\mathrm{ext},i},\\
z_j,&j\ne i.
\end{cases}
\]
Then
\[
z_{\mathrm{boundary}}\in
\mathcal B_{\mathrm{ext}}\setminus\mathcal B_{\mathrm{ext}}^\circ,
\qquad
\|z-z_{\mathrm{boundary}}\|_\infty
\le |z_i-z_i'|
\le\|z-z'\|_\infty.
\]
\Cref{proof:global-holder-extension:cutoff} makes the top derivative zero at both \(z'\) and
\(z_{\mathrm{boundary}}\). Hence
\[
\begin{aligned}
&\|\mathrm D^m\mathscr G g_{\mathrm{in}}(z)
-\mathrm D^m\mathscr G g_{\mathrm{in}}(z')\|_{\mathrm{op}}\\
&\quad=
\|\mathrm D^m\mathscr G g_{\mathrm{in}}(z)
-\mathrm D^m\mathscr G g_{\mathrm{in}}(z_{\mathrm{boundary}})\|_{\mathrm{op}}\\
&\quad\le
K_{\mathrm{cut},H}L_{\mathrm{in}}
\|z-z_{\mathrm{boundary}}\|_\infty^s
\le K_{\mathrm{cut},H}L_{\mathrm{in}}\|z-z'\|_\infty^s.
\end{aligned}
\]
Reversing the points treats the opposite crossing order. If both points lie outside the closed box, both derivatives are zero. Together with the closed-box estimate, these cases establish the displayed modulus for all \(z,z'\in\mathbb R^d\). In dimension zero the outside cases are empty.
% lean: Causalean.Mathlib.Analysis.Calculus.CubeExtension.exists_rectCutoffExtension_modulus_constant

\item \emph{Composition of the two radius factors.}
Define the fixed constants
\[
K_{\mathrm{cut}}:=\max\{K_{\mathrm{cut},D},K_{\mathrm{cut},H}\},
\qquad
K_{\mathrm{ext}}:=K_{\mathrm{cut}}K_{\mathrm{neigh}}>0.
\]
The box, cutoff, and all constants depend only on \(d,m,s\). For the given \(g,L\), choose the neighborhood response from \cref{proof:global-holder-extension:larger-box} and set
\[
L_{\mathrm{in}}:=K_{\mathrm{neigh}}L\ge0,
\qquad
g_{\mathrm{ext}}:=\mathscr G g_{\mathrm{neigh}}.
\]
Its agreement on the cube is
\[
g_{\mathrm{ext}}(z)=g_{\mathrm{neigh}}(z)=g(z),
\qquad z\in\mathcal C,
\]
and \cref{proof:global-holder-extension:cutoff} gives \(g_{\mathrm{ext}}\in C^m(\mathbb R^d)\). \Cref{proof:global-holder-extension:cutoff-derivatives,proof:global-holder-extension:global-modulus} give
\[
\|\mathrm D^k g_{\mathrm{ext}}(z)\|_{\mathrm{op}}
\le K_{\mathrm{cut}}(K_{\mathrm{neigh}}L)
=K_{\mathrm{ext}}L,
\qquad 0\le k\le m,
\]
and
\[
\|\mathrm D^m g_{\mathrm{ext}}(z)-\mathrm D^m g_{\mathrm{ext}}(z')\|_{\mathrm{op}}
\le K_{\mathrm{cut}}(K_{\mathrm{neigh}}L)\|z-z'\|_\infty^s
=K_{\mathrm{ext}}L\|z-z'\|_\infty^s.
\]
Every estimate uses multiplication by nonnegative radii, so the construction and bounds also apply when \(L=0\).
% lean: Causalean.Mathlib.Analysis.Calculus.CubeExtension.exists_global_holder_extension
\end{enumerate}
\end{proof}
